\chapter{O Sacramento do Matrimônio}

``Maridos, amai vossas mulheres, como Cristo amou a Igreja e se entregou por ela... Este mistério é grande; digo-o em referência a Cristo e à Igreja'' (Ef 5, 25.32). Com estas palavras do Apóstolo Paulo, o Catecismo da Igreja Católica enquadra toda a sua doutrina sobre o sacramento do Matrimônio, nos parágrafos 1601 a 1666. O Matrimônio não é uma invenção humana que a fé depois abençoa; é uma realidade que vem do Criador, ferida pelo pecado, elevada por Cristo e transformada em sinal eficaz da aliança de Deus com a humanidade. Cada casal cristão é convocado a ser, no coração do mundo, uma imagem viva daquele amor de Cristo pela Igreja: total, fiel, fecundo e irrevogável.

\section{O matrimônio no plano de Deus (1601--1620)}

A definição que o Catecismo apresenta de saída é precisa: ``A aliança matrimonial pela qual o homem e a mulher constituem entre si uma íntima comunhão de vida e de amor, foi fundada e dotada de suas próprias leis pelo Criador. Por sua natureza, é ordenada ao bem dos cônjuges e à procriação e educação dos filhos. Entre os batizados, foi elevada por Cristo Senhor à dignidade de sacramento'' (CDC, can. 1055 § 1). Três elementos articulam esta definição: a instituição divina do Matrimônio, a sua finalidade dupla --- o bem dos cônjuges e a abertura à vida --- e a sua elevação sacramental em Cristo.

A Sagrada Escritura começa com a criação do homem e da mulher à imagem e semelhança de Deus e termina com a visão das ``bodas do Cordeiro'' (Ap 19, 7-9). Entre esses dois polos, toda a Escritura fala do Matrimônio e do seu mistério, da sua instituição e do sentido que Deus lhe deu, da sua origem e do seu fim, das suas diversas realizações ao longo da história da salvação, das dificuldades nascidas do pecado e da sua renovação ``no Senhor'' na Nova Aliança de Cristo e da Igreja. A vocação ao Matrimônio está inscrita na própria natureza do homem e da mulher saídos das mãos do Criador; não é uma instituição puramente humana, apesar de variar nas diversas culturas. ``O bem da pessoa e da sociedade humana e cristã está estreitamente ligado à situação saudável da comunidade conjugal e familiar.''

Deus, que criou o homem por amor, chama-o ao amor: esta é a vocação fundamental e inata de todo ser humano. O homem foi criado à imagem e semelhança de Deus que é, ele mesmo, amor. Criando-o homem e mulher, Deus fez do amor mútuo entre eles uma imagem do amor absoluto e inesgotável com que Deus ama o homem. Este amor é bom --- ``muito bom'' aos olhos do Criador ---, e é chamado a ser fecundo: ``Sede fecundos, multiplicai-vos, enchei a terra'' (Gn 1, 28). A Escritura, além disso, é clara quanto à diferença e complementaridade do homem e da mulher: ``Não é bom que o homem esteja só'' (Gn 2, 18). A mulher, ``osso dos seus ossos e carne da sua carne'', é-lhe dada por Deus como companheira, igual na dignidade, semelhante em natureza. ``Por isso o homem deixará pai e mãe e se unirá à sua mulher, e os dois serão uma só carne'' (Gn 2, 24). O próprio Senhor mostrará que esta fórmula designa uma união indissolúvel.

Mas todo ser humano experimenta o mal em si mesmo e à sua volta. A desordem que se manifesta nas relações entre homem e mulher --- discórdia, espírito de dominação, infidelidade, ciúmes, conflitos que podem escalar ao ódio e à separação --- não deriva da natureza do homem e da mulher, nem da natureza das suas relações: deriva do pecado. O primeiro pecado, como ruptura com Deus, teve como primeira consequência a ruptura da comunhão originária entre o homem e a mulher: as acusações mútuas (cf.~Gn 3, 12), a atração recíproca --- dom do Criador --- transformada em relação de dominação e de concupiscência (cf.~Gn 3, 16b), e a vocação bela à fecundidade gravada pela dor do parto e pelo fardo do trabalho. No entanto, a ordem da criação persiste, mesmo que gravemente perturbada. Para curar as feridas do pecado, o homem e a mulher precisam da graça que Deus, em sua misericórdia infinita, nunca lhes recusa.

A pedagogia divina ao longo da história da salvação foi gradual. O Matrimônio ajudou o homem a superar o fechamento em si mesmo, o egoísmo, a busca do próprio prazer e a abrir-se ao outro, à ajuda mútua e à doação de si. A consciência moral acerca da unidade e indissolubilidade do Matrimônio desenvolveu-se sob a pedagogia da lei antiga. A poligamia dos patriarcas e reis não é ainda rejeitada explicitamente; mas a lei mosaica visava proteger a esposa da dominação arbitrária do marido. Os profetas, vendo na aliança de Deus com Israel uma imagem do amor conjugal exclusivo e fiel, foram preparando a consciência do Povo eleito para uma compreensão aprofundada da unidade e indissolubilidade do Matrimônio. Os livros de Rute e de Tobias testemunham de modo tocante um sentido elevado do Matrimônio e da fidelidade e ternura dos esposos. O Cântico dos Cânticos é, na leitura constante da Tradição, uma expressão única do amor humano, reflexo puro do amor de Deus --- amor ``forte como a morte'' que ``muitas águas não podem apagar'' (Ct 8, 6-7).

A aliança nupcial entre Deus e o Povo de Israel havia preparado o caminho para a aliança nova e eterna em que o Filho de Deus, encarnando-se e dando a sua vida, uniu a si de certa maneira toda a humanidade por ele salva. No limiar da sua vida pública, Jesus realiza seu primeiro sinal --- a pedido de sua Mãe --- durante uma festa de casamento em Caná (cf.~Jo 2, 1-11). A Igreja vê nesse sinal a confirmação da bondade do Matrimônio e o anúncio de que, a partir de então, o Matrimônio será um sinal eficaz da presença de Cristo. Na sua pregação, Jesus ensinou sem equívoco o sentido original da união do homem e da mulher como o Criador a quis desde o princípio: a permissão dada por Moisés de repudiar a esposa era uma concessão à dureza dos corações. ``O que Deus uniu, o homem não separe'' (Mt 19, 6). Este ensinamento pode parecer um fardo impossível de suportar; mas Jesus não impôs aos esposos um peso que não possam carregar. Vindo restaurar a ordem original da criação perturbada pelo pecado, Ele mesmo dá a força e a graça de viver o Matrimônio na nova dimensão do Reino de Deus. ``Seguindo Cristo, renunciando a si mesmos e carregando suas cruzes, os esposos poderão `receber' o sentido original do Matrimônio e vivê-lo com a ajuda de Cristo'' (cf.~Mt 19, 11).

É o que o apóstolo Paulo explicita: ``Maridos, amai vossas mulheres, como Cristo amou a Igreja e se entregou por ela para a santificar... Por isso o homem deixará pai e mãe e se unirá à sua mulher, e os dois serão uma só carne. Este mistério é grande; digo-o em referência a Cristo e à Igreja'' (Ef 5, 25-26.31-32). A vida cristã toda leva a marca do amor esponsal de Cristo e da Igreja. O próprio Batismo é um mistério nupcial: é o banho nupcial que precede a festa de bodas, a Eucaristia. O Matrimônio cristão torna-se, por sua vez, um sinal eficaz, o sacramento da aliança de Cristo e da Igreja.

Esta elevação sacramental do Matrimônio não anula nem diminui a vocação à virgindade consagrada; pelo contrário, as duas vocações se iluminam mutuamente e se fortalecem. Cristo é o centro de toda a vida cristã, e o vínculo com Ele tem precedência sobre qualquer outro vínculo familiar ou social. Desde o início da Igreja, houve homens e mulheres que renunciaram ao grande bem do Matrimônio para seguir o Cordeiro em toda parte (cf.~Ap 14, 4), para se ocupar exclusivamente das coisas do Senhor (cf.~1 Cor 7, 32) e para ir ao encontro do Esposo que vem. A virgindade pelo Reino é um desdobramento da graça batismal, um sinal poderoso da supremacia do vínculo com Cristo e da ardente esperança do seu retorno --- sinal que lembra também que o Matrimônio é uma realidade desta era presente que passa (cf.~Mc 12, 25; 1 Cor 7, 31). Tanto o sacramento do Matrimônio quanto a virgindade pelo Reino vêm do próprio Senhor; e, como diz São João Crisóstomo, quem deprecia o Matrimônio também diminui a glória da virgindade, pois o que parece bom somente em comparação com o mal não seria verdadeiramente bom.

\section{A celebração do sacramento (1621--1624)}

No Rito Latino, o Matrimônio entre dois fiéis católicos é normalmente celebrado durante a santa Missa, por causa da conexão de todos os sacramentos com o mistério pascal de Cristo. Na Eucaristia, memorial da Nova Aliança, é que Cristo se faz presente, e por isso convém que os esposos selarem o seu consentimento de se darem um ao outro à entrega das próprias vidas, unindo esta oferta à de Cristo à sua Igreja, tornada presente no sacrifício eucarístico, e que recebam a Eucaristia para, comungando num mesmo Corpo e num mesmo Cálice, ``formarem um só corpo'' em Cristo.

Porque a celebração litúrgica do Matrimônio é um ato sacramental de santificação, ela deve ser, em si mesma, válida, digna e fecunda. Por isso é conveniente que os noivos se preparem para a celebração recebendo o sacramento da Penitência.

Quanto ao ministro do sacramento, a Igreja latina ensina que os próprios esposos, como ministros da graça de Cristo, conferem mutuamente o sacramento do Matrimônio, expressando diante da Igreja o seu consentimento. Nas liturgias orientais, designadas pelo belo nome de ``Coroação'', o sacerdote ou bispo é ministro do sacramento: após receber o consentimento mútuo dos esposos, coroa-os em sinal da aliança matrimonial. As diversas tradições litúrgicas concordam, porém, em que a celebração inclui uma \emph{epiclese}, pela qual os esposos recebem o Espírito como comunhão de amor de Cristo e da Igreja, selo de sua aliança e fonte de seu amor e da força para permanecerem fiéis.

\section{O consentimento matrimonial (1625--1637)}

As partes do contrato matrimonial são um homem batizado e uma mulher batizada, livres para contrair Matrimônio, que expressam livremente o seu consentimento. ``Ser livre'' significa: não estar sob coação e não estar impedido por nenhuma lei natural ou eclesiástica. A Igreja considera que o mútuo consentimento dos esposos é o elemento indispensável que ``faz o matrimônio'' (CDC, can. 1057 § 1). Sem o consentimento, não há matrimônio. O consentimento é ``um ato humano pelo qual os cônjuges se doam e se recebem mutuamente'': ``Eu te recebo como minha esposa/meu esposo.'' Este consentimento, que une os esposos entre si, encontra o seu cumprimento nos dois ``tornando-se uma só carne'' (Gn 2, 24). O consentimento deve ser um ato da vontade de cada um dos contraentes, livre de coação ou de medo grave externo. Se esta liberdade falta, o Matrimônio é inválido. Por isso, a Igreja pode, após exame pela competente autoridade eclesiástica, declarar a nulidade de um matrimônio --- ou seja, que o matrimônio nunca existiu. Nesse caso, os contraentes ficam livres para casar, desde que cumpridas as obrigações naturais de uma eventual união anterior.

O sacerdote ou diácono que assiste à celebração do Matrimônio recebe o consentimento dos esposos em nome da Igreja e dá a bênção desta. A presença do ministro da Igreja --- e igualmente das testemunhas --- exprime visivelmente que o Matrimônio é uma realidade eclesial. Por isso a Igreja normalmente exige que os fiéis contraiam Matrimônio segundo a forma eclesiástica, por várias razões convergentes: o Matrimônio sacramental é um ato litúrgico e deve ser celebrado na liturgia pública da Igreja; ele introduz numa ordem eclesial e cria direitos e deveres na Igreja entre os cônjuges e em relação aos filhos; dado que o Matrimônio é um estado de vida na Igreja, é necessário ter certeza a seu respeito (daí a obrigação de ter testemunhas); e o caráter público do consentimento salvaguarda o ``sim'' dado e ajuda os esposos a permanecerem-lhe fiéis.

Para que o ``sim'' dos esposos seja um ato livre e responsável e para que a aliança matrimonial tenha fundamentos humanos e cristãos sólidos e duráveis, a preparação ao Matrimônio é de primordial importância. O exemplo e o ensino dados pelos pais e pelas famílias continuam sendo a forma especial desta preparação, mas o papel dos pastores e da comunidade cristã como ``família de Deus'' é indispensável para transmissão dos valores humanos e cristãos do Matrimônio e da família --- ``é absolutamente necessário instruir os jovens, a seu tempo, de preferência no seio da própria família, sobre a dignidade do amor conjugal, seu papel e seu exercício, para que, formados no respeito da castidade, possam, na idade conveniente, entrar num noivado honesto e contrair o próprio matrimônio'' (GS 49 § 3).

Nos países onde são frequentes os casamentos mistos --- entre um católico e um cristão não católico ---, estes requerem atenção particular da parte dos esposos e dos seus pastores. Os casamentos com disparidade de culto (entre um católico e uma pessoa não batizada) exigem ainda maior circunspecção. A diferença de confissão não constitui um obstáculo intransponível para o Matrimônio, quando os esposos conseguem colocar em comum o que cada um recebeu de sua comunidade de origem e aprendem mutuamente como cada um vive na fidelidade a Cristo. Entretanto, as dificuldades dos casamentos mistos não devem ser subestimadas: a separação dos cristãos ainda não foi superada, e os esposos arriscam-se a experimentar no seio da própria família a tragédia da divisão cristã. A disparidade de culto pode agravar essas dificuldades. Para os casamentos mistos, exige-se permissão da autoridade eclesiástica; para os casamentos com disparidade de culto, uma dispensa explícita é requerida para a validade. Em ambos os casos, os fiéis são chamados a conhecer e não excluir os fins e as propriedades essenciais do Matrimônio e as obrigações assumidas pelo cônjuge católico quanto ao batismo e à educação cristã dos filhos. O diálogo ecumênico procurou desenvolver, em muitas regiões, uma pastoral comum para os casamentos mistos, a fim de ajudar esses casais a viver a sua situação à luz da fé e a fazer florescer o que os une no respeito ao que os separa. Nos casamentos com disparidade de culto, o cônjuge católico tem uma tarefa particular: ``pois o marido não crente é santificado pela mulher, e a mulher não crente é santificada pelo marido crente'' (1 Cor 7, 14); e o amor conjugal sincero, a prática humilde e paciente das virtudes familiares e a perseverança na oração podem preparar o cônjuge não crente para aceitar a graça da conversão.

\section{Os efeitos do sacramento: aliança, vínculo e graça (1638--1654)}

Do Matrimônio válido nasce entre os esposos um vínculo que, por sua natureza, é perpétuo e exclusivo. Além disso, no Matrimônio cristão, os esposos são fortalecidos e como que consagrados por um sacramento especial, em vista dos deveres e da dignidade do seu estado (cf.~CDC, can. 1134). O consentimento pelo qual os esposos se doam e se recebem mutuamente é selado pelo próprio Deus: da aliança dos esposos nasce ``uma instituição confirmada pela lei divina, mesmo aos olhos da sociedade'' (GS 48 § 1). A aliança dos esposos é integrada na aliança de Deus com o homem: ``o amor conjugal autêntico é assumido no amor divino'' (GS 48 § 2). Assim, o vínculo matrimonial foi estabelecido pelo próprio Deus de tal maneira que um Matrimônio realizado e consumado entre batizados nunca pode ser dissolvido. Este vínculo, que resulta do ato humano livre dos esposos e da consumação do seu Matrimônio, é uma realidade doravante irrevogável, que dá origem a uma aliança garantida pela fidelidade de Deus. A Igreja não tem o poder de se opor a esta disposição da sabedoria divina.

A graça própria do sacramento do Matrimônio é destinada a aperfeiçoar o amor dos cônjuges e a robustecer a sua unidade indissolúvel. Por esta graça, eles ``se auxiliam mutuamente a santificar-se na vida conjugal e no acolhimento e na educação dos filhos'' (LG 11 § 2). Cristo é a fonte desta graça: como outrora Deus foi ao encontro do seu povo com uma aliança de amor e fidelidade, assim agora o nosso Salvador, Esposo da Igreja, vem ao encontro dos esposos cristãos pelo sacramento do Matrimônio. Ele habita com eles, dá-lhes a força de se carregarem mutuamente nas suas cruzes, de se perdoarem, de se ``sujeitar um ao outro no temor de Cristo'' (Ef 5, 21) e de se amarem com amor sobrenatural, terno e fecundo. Nas alegrias do amor e da vida familiar, Ele lhes dá já neste mundo um antegosto da festa de bodas do Cordeiro.

O amor conjugal envolve uma totalidade --- ``todo o conjunto dos elementos da pessoa entra nele: apelo do corpo e do instinto, força do sentimento e da afetividade, aspiração do espírito e da vontade'' (FC 13) ---: aponta para uma unidade profundamente pessoal que, para além da união na mesma carne, tende a formar um só coração e uma só alma; exige a indissolubilidade e a fidelidade na doação mútua definitiva; e é aberto à fecundidade. A unidade do Matrimônio exige, por sua própria natureza, a unicidade e a indissolubilidade da comunhão de pessoas entre os cônjuges, abraçando toda a sua vida: ``já não são dois, mas uma só carne'' (Mt 19, 6). São chamados a crescer continuamente na sua comunhão ``pela fidelidade cotidiana à promessa matrimonial de doação recíproca total'' (FC 19). Esta comunhão humana é confirmada, purificada e complementada pela comunhão em Jesus Cristo, dada pelo sacramento do Matrimônio. A fidelidade conjugal, por seu lado, tem a sua razão mais profunda na fidelidade de Deus à sua aliança e na de Cristo à sua Igreja: pelo sacramento do Matrimônio, os esposos são habilitados a representar e testemunhar essa fidelidade.

É verdade que pode parecer difícil, ou mesmo impossível, vincular-se a uma pessoa para toda a vida. Por isso, é tanto mais importante anunciar a Boa Nova de que a indissolubilidade do Matrimônio possui um sentido positivo, que a graça de Cristo torna realmente possível. Quando, contudo, a convivência se torna praticamente impossível por razões sérias, a Igreja admite a separação física dos cônjuges e a cessação da coabitação. Mas nesse caso os esposos não deixam de ser esposos diante de Deus; não ficam livres para contrair nova união enquanto ambos vivem.

Hoje, muitos fiéis recorrem ao divórcio civil e contraem novas uniões civis. A Igreja, por fidelidade às palavras de Jesus Cristo, mantém que não pode reconhecer como válida uma nova união, se o primeiro Matrimônio o era. Quem se encontra em tal situação, se os divorciados recasados não viveram como irmão e irmã, não podem aceder à comunhão eucarística. Não estão porém separados da Igreja: são chamados a participar da vida eclesial pela escuta da Palavra de Deus, pela frequência à Missa, pela perseverança na oração, pelas obras de caridade e pelas iniciativas da comunidade em favor da justiça, especialmente pela educação cristã dos seus filhos. A Igreja deve manifestar-lhes atenção solícita e amor acolhedor, para que eles não se sintam separados dela.

A abertura à fecundidade é uma das propriedades essenciais do Matrimônio. A fecundidade conjugal não se reduz à procriação: os cônjuges sem filhos podem ter uma vida conjugal plena de sentido, humana e cristã, dando a si mesmos por meio da adoção, por obras de caridade educativa ou por outras formas de serviço à vida. Mas o ``dom supremo do Matrimônio'' é a criança --- a vida humana nova que o amor dos esposos acolhe como dom de Deus e encarrega de educar. ``Os filhos são o dom mais excelente do Matrimônio e contribuem muito para o bem dos próprios pais'' (GS 50 § 1). A transmissão da vida é uma participação na obra criadora de Deus; e a educação das crianças na fé é uma das mais altas responsabilidades dos pais cristãos.

\section{A Igreja doméstica (1655--1658)}

Cristo escolheu nascer e crescer no lar da sagrada família de José e Maria. A Igreja não é outra coisa senão ``a família de Deus''. Nos nossos dias, em especial, as famílias que creem são de importância primordial como centros de fé viva e irradiante. E por isso o Concílio Vaticano~II retomou a bela expressão patrística: a família cristã é chamada de \textit{Ecclesia domestica}, Igreja doméstica.

No lar da família cristã é que o pai, a mãe, os filhos e todos os membros exercem o sacerdócio dos batizados: pela recepção dos sacramentos, pela oração e pela ação de graças, pelo testemunho de vida santa, pela abnegação e pelo amor ativo. O lar é, assim, a primeira escola de vida cristã e ``uma escola de riqueza humana'' (GS 52 § 1). É aqui que se aprende a trabalhar e a alegria do trabalho, o amor fraterno, o perdão generoso e até permanente, sobretudo o culto a Deus pela oração e pela oblação da própria vida.

A família cristã é uma comunhão de pessoas, sinal e imagem da comunhão do Pai e do Filho no Espírito Santo. Na geração e educação dos filhos ela reflete a obra criadora do Pai; pela mútua santificação dos cônjuges e pela educação dos filhos ela é um sinal da solicitude salvífica de Cristo pela Igreja, sua Esposa; e pela participação na oração e no sacrifício de Cristo ela é um reflexo do sacerdócio comum dos fiéis.

É preciso lembrar também o grande número de pessoas solteiras que --- muitas vezes em contra a suas vontades --- permanecem sozinhas. Estas pessoas têm frequentemente necessidade especial de atenção, bondade e dedicação. São Jesus no mundo, que permanecem abertas generosamente aos outros. Cristo crucificado é a fonte de todo amor, inclusive do amor a essas pessoas que não encontraram um caminho de vida estável. Elas são especialmente próximas do Coração de Jesus, e merecem por isso a afeição especial e a solicitude ativa da Igreja. As portas dos lares --- Igrejas domésticas --- e da grande família que é a Igreja devem estar abertas para todas elas.

\section{Em síntese (1659--1666)}

``Maridos, amai vossas mulheres, como Cristo amou a Igreja... Este mistério é grande; digo-o em referência a Cristo e à Igreja'' (Ef 5, 25.32). Com estas palavras do Apóstolo, o Catecismo abre a sua síntese sobre o sacramento do Matrimônio: tudo o que foi dito sobre a aliança matrimonial tem em Cristo e na Igreja o seu fundamento último e o seu modelo insuperável.

A aliança matrimonial pela qual o homem e a mulher constituem entre si uma parceria para toda a vida foi fundada e dotada de suas próprias leis pelo Criador. Por sua natureza, é ordenada ao bem dos cônjuges e à procriação e educação dos filhos. Entre batizados, foi elevada por Cristo à dignidade de sacramento.

O sacramento do Matrimônio significa a união de Cristo e da Igreja. Dá aos esposos a graça de se amarem com o amor de Cristo, aperfeiçoa o amor humano, robustece a unidade indissolúvel e santifica-os no caminho para a vida eterna.

O Matrimônio funda-se no consentimento dos contraentes --- na vontade de se darem mutuamente e de maneira definitiva, para viver uma aliança de amor fiel e fecundo. Como o Matrimônio estabelece o casal num estado público de vida na Igreja, é conveniente que a sua celebração seja pública, no seio de uma celebração litúrgica, diante de um sacerdote (ou de uma testemunha qualificada pela Igreja), das testemunhas e da assembleia dos fiéis.

A unidade, a indissolubilidade e a abertura à fecundidade são essenciais ao Matrimônio. A poligamia é incompatível com a unidade conjugal; o divórcio separa o que Deus uniu; a recusa da fecundidade priva o Matrimônio do seu ``dom supremo'', que é a criança.

Os divorciados que contraem uma nova união civil não estão separados da Igreja, mas não podem receber a comunhão eucarística; são, todavia, chamados a conduzir uma vida cristã, em especial pela educação dos seus filhos na fé.

O lar cristão é o lugar onde os filhos recebem o primeiro anúncio da fé. Por isso a família é chamada com razão ``Igreja doméstica'': ela é uma comunidade de graça e de oração, uma escola das virtudes humanas e da caridade cristã.
